% =====================================================================
%  Annexe A — Guide de soutenance
% =====================================================================
\chapter{Guide de soutenance}
\label{ann:soutenance}

Cette annexe rassemble ce qui est nécessaire le jour de la présentation : les comptes à
utiliser, un déroulé de démonstration minuté, et un jeu de questions probables avec des
réponses appuyées sur le code réel.

% =====================================================================
\section{Préparation}

\subsection{Lancement de l'environnement}

Trois terminaux, dans cet ordre.

\begin{enumerate}
  \item \textbf{Base de données.} Vérifier que PostgreSQL est démarré et que la base
    \code{medicacom} existe.
  \item \textbf{Backend.} Depuis \code{backend/}, activer l'environnement virtuel puis
    lancer le serveur sur le port~8000. Sous PowerShell, l'activation se fait avec
    \code{.\textbackslash.venv\textbackslash Scripts\textbackslash Activate.ps1} ; sous
    Git~Bash, avec \code{source .venv/Scripts/activate}. La confusion entre les deux
    syntaxes est la cause la plus fréquente de l'erreur
    « \code{uvicorn} n'est pas reconnu ».
  \item \textbf{Frontend.} Depuis \code{frontend/}, \code{npm run dev}. L'interface
    répond sur le port~5173.
\end{enumerate}

Ouvrir un onglet supplémentaire sur \code{http://localhost:8000/docs}. La documentation
interactive de l'API est utile si le jury demande une preuve côté serveur.

\subsection{Comptes de démonstration}

Tous les comptes ci-dessous partagent le mot de passe défini par la variable
\code{DEMO\_PASSWORD}. Ils sont réservés au développement et à la démonstration.

\begin{table}[H]
\centering
\caption{Comptes du jeu de démonstration}
\label{tab:comptes}
\small
\begin{tabular}{@{}L{5.3cm} L{3.4cm} L{5.1cm}@{}}
\toprule
\textbf{Adresse} & \textbf{Rôle} & \textbf{Écran d'accueil} \\
\midrule
admin@medicacom.com            & Admin. Medicacom & Administration \\
admin.entreprise@medicacom.com & Admin. entreprise & Administration \\
marwen.b@medicacom.com         & Direction Générale & Tableau de bord Direction \\
l.trabelsi@medicacom.com       & Resp. département & Tableau de bord Département \\
h.gharbi@medicacom.com         & Resp. département & Tableau de bord Département \\
c.lambert@medicacom.com        & Resp. département & Tableau de bord Département \\
a.rousseau@medicacom.com       & Contributeur métier & Opportunités \\
k.zouari@medicacom.com         & DSI / Architecte & Architecture \& données \\
n.hamdi@medicacom.com          & Data Office & Architecture \& données \\
y.benali@medicacom.com         & AI Lead & Portefeuille \\
m.dupuis@medicacom.com         & Finance & Valeur \& ROI \\
f.moreau@medicacom.com         & PMO & Portefeuille \\
audit@medicacom.com            & Auditeur & Portefeuille \\
\bottomrule
\end{tabular}
\end{table}

% =====================================================================
\section{Déroulé de démonstration (10 à 15 minutes)}

\subsection*{Séquence 1 — Connexion et parcours (1 min 30)}

\textbf{Où :} page de connexion.\quad\textbf{Quoi :} se connecter avec
\code{marwen.b@medicacom.com}.

\begin{oral}
Medicacom AOI Transformation répond à un problème simple à énoncer et difficile à
résoudre : dans un laboratoire pharmaceutique, les idées d'IA ne manquent pas, mais rien
ne permet de décider lesquelles financer. La plateforme couvre le parcours complet, du
diagnostic de maturité à la mesure de la valeur réellement captée. Je me connecte en
Direction Générale.
\end{oral}

\subsection*{Séquence 2 — Tableau de bord Direction (1 min 30)}

\textbf{Où :} écran d'accueil.\quad\textbf{Quoi :} montrer les quatre indicateurs, puis
la matrice Impact~$\times$~Faisabilité.

\begin{oral}
Ces quatre chiffres ne sont pas saisis. Ils sont agrégés à la demande depuis les
projets, les coûts et les bénéfices. Le ROI de portefeuille rapporte les gains cumulés
sur trois ans au coût total de possession. Dans la matrice, chaque bulle est une
opportunité ; sa taille est le bénéfice annuel attendu, et le quadrant supérieur droit
est ce que l'organisation devrait lancer en premier.
\end{oral}

\subsection*{Séquence 3 — Diagnostic et readiness (2 min)}

\textbf{Où :} menu \emph{Diagnostic}.\quad\textbf{Quoi :} montrer la jauge, les deux
bandeaux d'alerte, puis descendre jusqu'au tableau des preuves.

\begin{oral}
Le readiness est le point de départ. Chaque dimension est notée de 0 à 5, pondérée par
son importance et par la confiance que le répondant accorde à sa propre réponse. Cette
confiance intervient au numérateur et au dénominateur, donc une réponse mal étayée ne
fausse le score dans aucun sens. Les deux dimensions en ambre sont bloquantes : sous le
seuil et marquées critiques, elles interdisent un passage en production. Et chaque note
est adossée à une preuve — c'est ce qui rend le diagnostic auditable.
\end{oral}

\textbf{Action facultative :} ouvrir « Poursuivre le diagnostic », modifier une note,
enregistrer, montrer que le score bouge.

\subsection*{Séquence 4 — Opportunité et qualification (2 min 30)}

\textbf{Où :} menu \emph{Opportunités}, puis « Nouvelle opportunité ».\quad
\textbf{Quoi :} d'abord signaler le bandeau de doublons, puis créer une opportunité.

\begin{oral}
Le bandeau ambre est un apport du croisement interservices : deux départements ont
exprimé un besoin très proche sans le savoir, et la plateforme les rapproche sur leurs
termes partagés. Sans cet écran, les deux projets seraient financés séparément.

Je crée maintenant une opportunité. Ce questionnaire n'est pas décoratif : chaque
réponse est un fait d'entrée du moteur. Je déclare des données disponibles mais non
structurées, du texte libre à traiter, pas de génération de contenu. Je valide — et la
qualification s'exécute immédiatement. Il n'y a pas d'étape séparée : dans cette
plateforme, il n'existe pas d'opportunité non qualifiée.
\end{oral}

\subsection*{Séquence 5 — Portefeuille et repriorisation (2 min)}

\textbf{Où :} menu \emph{Portefeuille}, puis « Recalculer la priorisation ».

\begin{oral}
Le score de chaque ligne combine six critères pondérés, moins une pénalité de risque
proportionnelle au produit probabilité par impact. Le rang et la vague en découlent.

Si la Direction décide que le contexte impose de privilégier l'impact financier, elle
pousse ce curseur et relance le calcul. Je le fais : douze opportunités changent de
rang. Et ce changement de pondération est lui-même tracé dans le journal d'audit, avec
l'ancienne et la nouvelle grille.
\end{oral}

\subsection*{Séquence 6 — Fiche projet (2 min 30)}

\textbf{Où :} cliquer sur la première ligne du portefeuille.

\begin{oral}
C'est le livrable que reçoit un comité d'investissement. Trois choses le distinguent
d'une présentation classique. D'abord, la recommandation est justifiée : les pastilles
sont les règles satisfaites, et l'alternative écartée est nommée avec son motif.
Ensuite, le score dans la colonne de droite est décomposé critère par critère, pénalité
de risque comprise. Enfin, les garde-fous sont générés automatiquement : sur une
opportunité en périmètre GxP, la plateforme ajoute la validation informatique, la
traçabilité ALCOA+ et la revue humaine obligatoire.
\end{oral}

\subsection*{Séquence 7 — Changement de rôle (2 min)}

\textbf{Où :} se déconnecter, se reconnecter avec \code{l.trabelsi@medicacom.com}.

\begin{oral}
Même application, même URL d'accueil, écran différent. Le menu latéral a changé aussi :
il ne contient que les écrans auxquels ce rôle a droit.

Mais l'essentiel n'est pas ce qu'on voit. Masquer un bouton dans React n'est pas un
contrôle d'accès. Si j'appelais l'API directement pour lire une opportunité d'un autre
département, je recevrais 404 — et non 403, parce que répondre 403 confirmerait que la
ressource existe. Quarante tests vérifient ce comportement rôle par rôle, en appelant
l'API sans passer par l'interface.
\end{oral}

\subsection*{Séquence 8 — Valeur et audit (1 min 30)}

\textbf{Où :} menus \emph{Valeur \& ROI} puis \emph{Journal d'audit}.

\begin{oral}
L'écran Valeur ferme la boucle : le business case ne sert plus seulement à obtenir un
budget, il compare la valeur promise à la valeur mesurée.

Et le journal répond à la question « pourquoi cette décision ? » six mois plus tard.
Chaque modification laisse son auteur, sa date, l'ancienne et la nouvelle valeur. Point
important : il n'existe aucun moyen, ni par l'interface ni par l'API, de modifier ou de
supprimer une entrée. Le journal est en insertion seule, ce qui est une exigence directe
du contexte pharmaceutique.
\end{oral}

\subsection*{Conclusion orale (30 s)}

\begin{oral}
Un dernier point, qui est aussi le plus important. Le produit s'appelle « AOI
Transformation » et recommande des solutions d'intelligence artificielle. Il n'en
contient aucune. Le moteur qui produit ces recommandations est un système à base de
règles pondérées, écrites par un expert, entièrement déterministe. C'est un choix
assumé : sans historique de décisions, un modèle n'aurait rien à apprendre, et en
périmètre réglementé une décision non explicable est une décision indéfendable.
L'intégration d'un modèle de langage figure dans les perspectives du rapport, pas dans
les fonctionnalités livrées.
\end{oral}

% =====================================================================
\section{Questions probables du jury}

\subsection*{Sur l'intelligence artificielle}

\textbf{1. Quel modèle d'IA avez-vous utilisé ?}\\
Aucun. La plateforme ne contient ni modèle entraîné, ni appel à un fournisseur externe.
Ce que le nom du produit désigne, c'est le domaine traité — la gouvernance de l'adoption
de l'IA — pas la technologie employée. Le moteur qui produit les recommandations est un
système expert à base de 39~règles pondérées réparties sur 8~types de solution.

\textbf{2. Comment pouvez-vous l'affirmer ?}\\
Ni \code{requirements.txt} ni \code{package.json} ne déclarent de bibliothèque
d'apprentissage ou de client de fournisseur de modèle. Aucun fichier de modèle entraîné
n'existe dans le dépôt, aucune clé d'API n'est attendue par la configuration, et le
backend n'émet aucune requête HTTP sortante. Le seul client HTTP installé sert aux
tests.

\textbf{3. Pourquoi ne pas avoir intégré un modèle ?}\\
Pour quatre raisons. Il n'existe pas de données d'entraînement : prédire le bon type de
solution supposerait des centaines de besoins qualifiés puis validés par leur résultat
réel, corpus dont aucune organisation ne dispose au démarrage. La décision doit être
contestable règle par règle en comité. Le contexte GxP pénalise l'opacité. Et le
déterminisme rend le moteur testable sur des sorties exactes.

\textbf{4. Quelle est la différence entre votre moteur et de l'apprentissage
automatique ?}\\
L'origine de la connaissance. Ici, elle est écrite par un humain et lisible dans le
code ; en apprentissage, elle est inférée statistiquement de données. Il en découle tout
le reste : pas de jeu d'entraînement, déterminisme, explicabilité native — mais aussi
l'incapacité à découvrir ce qu'un expert n'a pas prévu.

\textbf{5. Que ferait un modèle de langage s'il était ajouté ?}\\
Trois choses, toutes en entrée du moteur, jamais à sa place : pré-remplir les quinze
faits de qualification à partir d'une description libre, rédiger un premier jet de fiche
projet, et détecter des doublons sémantiques là où la comparaison lexicale échoue. Le
modèle proposerait des entrées ; les règles continueraient de décider, pour préserver la
traçabilité.

\subsection*{Sur les algorithmes}

\textbf{6. Comment calculez-vous le readiness ?}\\
Chaque dimension porte une note sur 5, un poids et une confiance entre 0 et 1. Le score
est cent fois la somme des notes pondérées par poids et confiance, divisée par la somme
des maxima pondérés de la même façon. La confiance apparaît aux deux termes, donc une
réponse mal étayée s'efface au lieu de fausser le résultat.

\textbf{7. Pourquoi les moyennes d'axe n'utilisent-elles pas la confiance ?}\\
Parce qu'elles décrivent deux choses différentes. Le score global mesure une maturité
consolidée, pour laquelle la solidité du constat compte. La moyenne d'axe décrit le
niveau observé, indépendamment de la qualité de l'instruction.

\textbf{8. Que se passe-t-il si une dimension n'a pas été renseignée ?}\\
Une confiance nulle la marque comme non instruite. Elle sort du calcul du score et des
moyennes — une dimension non mesurée ne vaut pas zéro — et pèse uniquement sur le taux
de complétion. Sans ce mécanisme, ajouter une question au référentiel ferait chuter le
readiness de toutes les campagnes existantes sans qu'aucune réalité n'ait changé.

\textbf{9. Comment le score de priorité est-il construit ?}\\
Six critères notés sur 100, pondérés, moins une pénalité de risque. Les pondérations par
défaut sont 0,25 pour l'impact métier, 0,20 pour l'impact financier, 0,20 pour la
faisabilité, 0,15 pour l'alignement, 0,10 pour le readiness et 0,10 pour les synergies.
La pénalité vaut le produit probabilité par impact, ramené sur 25, multiplié par un
plafond de 15~points.

\textbf{10. Que se passe-t-il si l'utilisateur saisit des poids dont la somme n'est pas
1 ?}\\
Ils sont renormalisés avant application. C'est un choix d'ergonomie : l'utilisateur
pousse le critère qui l'intéresse sans avoir à rééquilibrer les cinq autres à la main,
et les scores restent comparables.

\textbf{11. Pourquoi 15 points de pénalité maximale, et pas 20 ou 10 ?}\\
C'est un étalonnage, pas une constante théorique. Assez élevé pour qu'une opportunité
très risquée soit réellement déclassée, assez bas pour qu'un risque maîtrisable
n'annule pas un impact métier fort. La valeur est isolée dans une constante nommée,
modifiable en un point.

\textbf{12. Pourquoi le payback n'utilise-t-il pas le TCO ?}\\
Parce qu'il répond à une question précise : en combien de mois l'investissement initial
est-il récupéré ? Y intégrer trois années de coûts récurrents non encore engagés
allongerait artificiellement le délai. Le numérateur est donc le CAPEX seul, le
dénominateur le gain annuel net des coûts d'exploitation, ramené au mois.

\textbf{13. Qu'est-ce qu'une règle bloquante dans le moteur de qualification ?}\\
Une condition sans laquelle le type de solution n'a pas de sens — par exemple des
données d'image pour de la vision par ordinateur. Si elle n'est pas satisfaite, le type
est disqualifié et son score divisé par quatre. Il reste visible dans les alternatives
avec le motif de sa disqualification : savoir qu'un type serait pertinent mais qu'il est
bloqué sur la vérité terrain est plus utile qu'une simple absence.

\subsection*{Sur la sécurité}

\textbf{14. Comment garantissez-vous qu'une organisation ne voit pas les données d'une
autre ?}\\
L'identifiant d'organisation est lu dans le jeton signé, jamais dans le corps ou les
paramètres de la requête. Un utilisateur ne peut donc pas se déclarer membre d'une autre
organisation sans forger le jeton, ce qui exigerait la clé de signature. Treize tests
construisent une seconde organisation complète et vérifient qu'aucune lecture, écriture,
suppression ou affectation ne traverse la frontière.

\textbf{15. Pourquoi 404 et non 403 pour une ressource hors périmètre ?}\\
Répondre 403 confirmerait que la ressource existe, ce qui est déjà une fuite. En
revanche, une écriture refusée \emph{à l'intérieur} de sa propre organisation répond
bien 403 : l'utilisateur voit l'objet, il n'a simplement pas le droit d'agir dessus.

\textbf{16. Si je masque un bouton dans React, suis-je protégé ?}\\
Non, et c'est un point que le projet prend au sérieux. Le masquage est une commodité
d'usage. La protection réelle est la dépendance d'autorisation côté serveur. Les tests
appellent les endpoints directement, sans passer par l'interface, précisément parce que
c'est ainsi qu'un attaquant procéderait.

\textbf{17. Un administrateur d'entreprise peut-il devenir administrateur de la
plateforme ?}\\
Non. La permission \code{platform.admin} est la seule à franchir la frontière entre
organisations, et une vérification explicite empêche un administrateur d'entreprise de
se l'attribuer ou de la placer dans un rôle qu'il crée. Deux tests couvrent ces deux
chemins.

\textbf{18. Comment les mots de passe sont-ils stockés ?}\\
Hachés avec bcrypt via passlib. Aucun mot de passe en clair n'est stocké ni journalisé.
Les comptes de démonstration partagent un mot de passe lu d'une variable
d'environnement, jamais écrit dans le code de l'interface.

\textbf{19. Qu'est-ce qui empêche de modifier le journal d'audit ?}\\
L'API n'expose aucune route de mise à jour ni de suppression sur cette table. La
traçabilité repose donc sur l'absence de chemin d'écriture, pas sur une permission qu'on
pourrait accorder par erreur. La table est aussi dépourvue de clés étrangères, pour
qu'une trace survive à la suppression de l'entité qu'elle décrit.

\textbf{20. Quelles faiblesses de sécurité connaissez-vous dans la version actuelle ?}\\
Trois, assumées. La clé de signature et le mot de passe du superutilisateur ont des
valeurs par défaut à redéfinir avant production. Il n'existe ni jeton de
rafraîchissement ni liste de révocation, donc un jeton volé reste exploitable jusqu'à
son expiration. Et il n'y a pas d'authentification à deux facteurs, ce qui serait
attendu en périmètre réglementé.

\subsection*{Sur l'architecture et le code}

\textbf{21. Pourquoi FastAPI plutôt qu'un autre framework ?}\\
Pour trois propriétés utiles ici : la validation des entrées adossée à Pydantic, la
documentation OpenAPI générée automatiquement — précieuse pour démontrer le
comportement de l'API — et le système d'injection de dépendances, qui a permis de
centraliser l'autorisation en deux fonctions plutôt que de la disperser dans 99~routes.

\textbf{22. Pourquoi séparer les services des routers ?}\\
Pour que les formules soient testables sans serveur ni base de données. C'est ce qui a
permis de corriger le calcul du payback et la sémantique du readiness tard dans le
projet sans redouter de régression.

\textbf{23. Comment gérez-vous l'évolution du schéma ?}\\
Par trois migrations Alembic versionnées, écrites pour s'appliquer sur une base déjà
peuplée. La troisième, qui a introduit les périmètres, crée les colonnes avec une valeur
par défaut le temps de la reprise, recopie les rattachements existants, puis retire la
valeur par défaut.

\textbf{24. Pourquoi la table utilisateur s'appelle-t-elle \code{app\_user} ?}\\
Parce que \code{user} est un mot réservé en PostgreSQL.

\textbf{25. Pourquoi conserver \code{department\_id} sur l'utilisateur alors qu'une
table d'association existe ?}\\
La colonne désigne le département principal et sert de repli ; la table d'association
porte les rattachements multiples introduits plus tard. Cette coexistence est le
résultat d'une migration qui devait préserver les données existantes, et non un choix de
conception initial.

\subsection*{Sur le périmètre livré}

\textbf{26. Tous les modules du cahier des charges sont-ils terminés ?}\\
Neuf sur douze le sont. Trois ne le sont que partiellement, et le rapport les détaille :
l'export de la fiche projet produit du HTML imprimable et non un fichier \code{.docx},
le suivi de projet dispose de son API mais pas de son écran dédié, et les notifications
sont persistées sans panneau de consultation.

\textbf{27. Le bouton dit « Export Word / PDF » — est-ce exact ?}\\
Non, et c'est signalé comme tel dans le rapport. L'API renvoie du JSON et une page HTML
mise en page pour l'impression. L'utilisateur obtient un PDF via son navigateur, mais
aucun fichier Word n'est généré. Le libellé est en avance sur la fonction.

\textbf{28. Qu'est-ce qui a été le plus difficile ?}\\
La reprise du contrôle d'accès. La première version ne comptait que dix permissions,
toutes en écriture : les quarante-six endpoints de lecture n'étaient protégés que par
l'authentification, ce qui laissait tout utilisateur connecté lire l'intégralité de son
organisation. Le durcissement a porté le catalogue à 41~permissions et introduit la
notion de périmètre, sans casser les modules existants.

\textbf{29. Comment avez-vous validé que le durcissement ne cassait rien ?}\\
Par une suite de 112~tests, dont 40 couvrent le comportement de chaque rôle et 4 sont
des garde-fous structurels qui parcourent la table de routage pour vérifier qu'aucun
endpoint n'est exposé sans contrôle et qu'aucune dépendance d'autorisation n'a été
silencieusement transformée en paramètre de requête.

\textbf{30. Que feriez-vous différemment en recommençant ?}\\
Je concevrais le modèle de permissions avant les fonctionnalités, et non après. Le
rattraper sur un code déjà écrit a coûté plus cher que de le poser d'emblée, et c'est
la difficulté la plus instructive du projet.
